
\documentclass{article}
\usepackage{amsmath}
\usepackage{amssymb}
\usepackage{geometry}
\usepackage{enumitem}
\geometry{margin=1in}

\title{Sequential Variable Elimination for Rational Inequalities}
\author{Enhanced Rational Inequality Solver}
\date{\today}

\begin{document}
\maketitle

\section{Introduction}

This document shows the step-by-step elimination of variables from a system of rational inequalities. 
For each variable elimination, we show four sub-steps:
\begin{enumerate}
\item Original inequalities involving the variable
\item Min/Max expressions after isolation
\item Min/Max expressions after variable elimination
\item Final expanded inequalities
\end{enumerate}

\section{Initial System}

Starting with the following system of inequalities:

\begin{align}
& \frac{1}{x_{1} + 2} >= 0 \\
& \frac{x_{2}}{x_{1} - 1} <= 3 \\
& x_{1} + x_{2} >= 0 \\
& x_{1} - 2 x_{2} <= 5 \\
& a + b <= c
\end{align}

\section{Step 1: Eliminate Variable $x1$}

\subsection{Step 1.1: Input Inequalities}

Inequalities involving $x1$:
\begin{align}
& \frac{1}{x_{1} + 2} >= 0 \\
& \frac{x_{2}}{x_{1} - 1} <= 3 \\
& x_{1} + x_{2} >= 0 \\
& x_{1} - 2 x_{2} <= 5
\end{align}

Inequalities not involving $x1$ (preserved):
\begin{align}
& a + b <= c
\end{align}

\subsection{Step 1.2: Min/Max Bounds for $x1$}

\begin{align}
& \max\left(\frac{x_{2}}{3} + 1, - x_{2}\right) \leq x1 \leq \min\left(\frac{x_{2}}{3} + 1, 2 x_{2} + 5\right)
\end{align}

\subsection{Step 1.3: Constraint After Eliminating $x1$}

For the system to be feasible:
\begin{align}
& \max\left(- x_{2}, \frac{x_{2}}{3} + 1\right) \leq \min\left(\frac{x_{2}}{3} + 1, 2 x_{2} + 5\right)
\end{align}

\subsection{Step 1.4: All Resulting Inequalities}

\begin{align}
& a + b <= c \\
& - x_{2} <= \frac{x_{2}}{3} + 1 \\
& - x_{2} <= 2 x_{2} + 5 \\
& \frac{x_{2}}{3} + 1 <= \frac{x_{2}}{3} + 1 \\
& \frac{x_{2}}{3} + 1 <= 2 x_{2} + 5
\end{align}

\section{Step 2: Eliminate Variable $x2$}

\subsection{Step 2.1: Input Inequalities}

Inequalities involving $x2$:
\begin{align}
& - x_{2} <= \frac{x_{2}}{3} + 1 \\
& - x_{2} <= 2 x_{2} + 5 \\
& \frac{x_{2}}{3} + 1 <= \frac{x_{2}}{3} + 1 \\
& \frac{x_{2}}{3} + 1 <= 2 x_{2} + 5
\end{align}

Inequalities not involving $x2$ (preserved):
\begin{align}
& a + b <= c
\end{align}

\subsection{Step 2.2: Min/Max Bounds for $x2$}

\begin{align}
& x2 \leq \min\left(- \frac{3}{4}, - \frac{5}{3}, - \frac{12}{5}\right)
\end{align}

\subsection{Step 2.3: Constraint After Eliminating $x2$}

Upper bound constraint:
\begin{align}
& \text{Any value} \leq - \frac{12}{5}
\end{align}

\subsection{Step 2.4: All Resulting Inequalities}

\begin{align}
& a + b <= c \\
& \frac{x_{2}}{3} + 1 <= \frac{x_{2}}{3} + 1
\end{align}


\section{Summary}

The sequential variable elimination process transforms the original system of inequalities 
into a new system where the specified variables have been eliminated. Each elimination 
step potentially increases the number of inequalities due to the expansion of min/max 
combinations, but removes one variable from consideration.

This approach is particularly useful for:
\begin{itemize}
\item Eliminating intermediate variables in optimization problems
\item Reducing dimensionality of constraint systems
\item Analyzing parameter dependencies in rational inequalities
\end{itemize}

\end{document}
